\section{Related Work}
\label{sec:related}

\para{Coulombic efficiency measurement and interpretation.}
CE is the primary metric for lithium inventory loss in rechargeable batteries. \citet{xiao2020understanding} established rigorous measurement protocols for Li metal cells, showing that CE alone does not predict cycle life---electrolyte volume, lithium thickness, and testing protocol all influence measured values. \citet{hobold2021beyond} provided a comprehensive perspective on CE determinants, identifying two regimes: Li$^0$-dominated loss below ${\sim}95\%$ CE and SEI-dominated loss above ${\sim}95\%$ CE. They emphasized that at the high-CE frontier ($>99.9\%$), trace impurities become critical but understudied variables. Our work directly addresses this gap by systematically modeling impurity effects in the SEI-dominated regime.

\para{SEI composition and performance.}
The composition of the SEI determines both CE and safety. \citet{hobold2024lithiumoxide} overturned the prevailing LiF-centric design paradigm by showing that Li$_2$O prevalence correlates more strongly with high CE across ten diverse electrolytes. They found Li$_2$O to be the most abundant identifiable SEI phase (14.6\% of capacity loss in standard carbonate electrolyte), and demonstrated that fluorine-free electrolytes can achieve $>99\%$ CE when Li$_2$O-rich SEI forms. \citet{otto2023sei} developed coulometric titration methods to quantify SEI growth on Li metal anodes in solid-state batteries. Our simulation approach models how impurities shift SEI composition toward either beneficial (Li$_2$O-rich, inorganic) or detrimental (organic-rich, resistive) phases.

\para{Impurity effects on battery performance.}
\citet{fink2021metallic} conducted the most comprehensive experimental study of metallic impurities, testing Fe, Al, Mg, Cu, and Si at 1~wt\% in both anodes and cathodes. They found that metallic contaminants at the anode disrupt performance through direct reaction with lithium and catalysis of electrolyte degradation, while cathode-side contaminants cross over during formation to disrupt SEI. However, they also noted that low-level doping can increase capacity retention and thermal stability. \citet{choe2024reevaluation} demonstrated that ${\sim}1$~mol\% Mg in Li$_2$CO$_3$ precursors improved NCM622 capacity retention from 76.1\% to 82.6\% at 200 cycles. Neutron diffraction revealed that $>80\%$ of Mg$^{2+}$ occupies lithium sites when introduced through the lithium carbonate route, providing structural stabilization. Unlike these single-impurity studies, we compare four impurity types within a unified simulation framework and map dose-response curves across the ppm range.

\para{Water impurities and thermal safety.}
\citet{baakes2023impact} developed a mechanistic thermal-runaway model with 12 degradation reactions and 20 species, showing that SEI thickness and composition dominate safety behavior more than water content alone. They found that water at 168--260~ppm shows no safety difference from the reference case, while thick inorganic SEI delays self-heating onset by ${\sim}19^\circ$C. Water plays a dual role: it acts as an endothermic heat sink through the PF$_5$ dissolution reaction while also driving exothermic LiOH decomposition. We incorporate their findings into our parameterization of water's concentration-dependent effects.

\para{Computational battery modeling.}
Physics-based battery models, particularly the Doyle--Fuller--Newman (DFN) family implemented in \pybamm~\citep{pybamm2020}, enable systematic parameter sweeps that would require months of physical experimentation. \citet{chen2020parameterization} provided the reference parameterization for NMC/graphite pouch cells that we adopt. \citet{single2021impedance} developed impedance spectroscopy theory for lithium batteries that informs our SEI resistivity modeling. While these models are well-validated for bulk electrochemical behavior, they have not been used to study trace impurity effects on SEI formation---a gap our work addresses.
